\section{What empirical evidence can establish}
\label{sec:empirical-interpretation}

The prior's support condition quantifies over every Boolean input and every
allowed simultaneous disturbance. A finite-noise certificate at a held-out
image proves stability at that image only. A successful adversarial search
provides a counterexample when both its disturbance budget and its score
change are verified. Failure to find an attack proves neither stability
nor membership in the prior. Floating-point sufficient tests are reported
separately from interval or exact-arithmetic checks.

There are two distinct empirical eligibility questions. First, does a
trained network satisfy the stability restriction? Second, does it satisfy
the likelihood, which requires every signed training output to exceed
$1/2$? A stable network that has training mistakes is evidence about the
condition's occurrence under training, but is not a draw admissible under
the noiseless posterior for that training set. Conversely, successful
optimization alone provides no stability guarantee. All declared runs,
including optimization failures, are retained.

The original all-input clean-margin convention imposes a further
requirement beyond training eligibility. The image counterexamples near
zero output already fail that clean-margin requirement at their witness
input. They refute automatic global stability of the trained
parameterization; they are not examples already inside that stronger
base support. In contrast, selected exhaustive Boolean-domain positives
verify absolute clean margins above $1/2$ everywhere, as well as all
signed training margins and internal stability. The separate readout
rescaling result explains how a finite-domain parameterization can enter
both supports, but charges the resulting Gaussian energy cost.

\paragraph{Finite grayscale images and Boolean inputs.}
An eight-bit pixel $p\in\{0,\ldots,255\}$ normalized as $x=2p/255-1$
can be written exactly as
\[
 x=\sum_{b=0}^{7}\frac{2^b}{255}B_b,\qquad B_b\in\{-1,1\}.
\]
Composing the first affine layer with this fixed linear encoding represents
the same image network as a Boolean-input tanh network with eight times
as many input coordinates. All hidden computations and internal disturbance
responses are unchanged at corresponding inputs. This observation connects
the stability property to a finite image domain; it does not equate the
independent Gaussian prior on pixel weights with the independent Gaussian
prior on bit weights, and it does not turn a test-set check into a proof
over all possible images.

\paragraph{Why a continuous-box counterexample has a different scope.}
When the input dimension exceeds the first hidden width, one may solve
$W^1\bar x=0$ and sometimes obtain $x\in[-1,1]^d$. With no later biases,
every clean hidden state and the readout then vanish. This exposes
amplification that saturation at ordinary data inputs can hide. Such an
input is a useful stress test for a continuous-box version of the condition.
It is a counterexample to the Boolean-domain condition only if it is a
Boolean input, or if a separate valid discrete input is found and verified.
Likewise, an interpolated image need not lie on the exact eight-bit grid.

\paragraph{Interpreting calibration.}
The factor $\sqrt{1+HE_v}$ is the worst-case noise gain for hidden layers
that are individually nonexpansive in Euclidean norm. Hence every such
network satisfies the score condition when $c\le1/16$. Smaller $c$ allows
greater amplification relative to this reference. A positive result only
at a small radius should not be described as stability to arbitrary
percent-level changes of every activation. Both the normalized total
budget and the readout calibration are part of the claim.

Global failures motivate studying fixed trusted input regions, chosen
independently of the labeled sample, on which the restriction is imposed.
If such a region has input-distribution mass at least $1-\rho$, the
statistical proof adds at most $\rho$ to the test error. Proving stability
throughout the region remains necessary; merely observing no failures on
a finite set of images does not supply that proof.
